\section{Related Work}
\label{sec:related}

Most crack analysis detects or segments cracks in individual images~\cite{zou2012cracktree,liu2019deepcrack,eisenbach2017gaps}.
Consolidated datasets~\cite{kulkarni2022crackseg9k,benz2024omnicrack30k} keep this image-level formulation without cross-observation identities.
Multi-view pipelines fuse images from one inspection into surface models or segmentations~\cite{benz2025crackstructures} but do not explicitly represent persistent identities across observations.
Structural health monitoring tracks known cracks over time with fixed cameras, designated locations~\cite{parente2022monitoring,harb2025multitemporal,sun2025perspective}, or guided setups~\cite{nec2019guide}, assigning identities when monitoring is set up; our task is open-set association.

Our setting combines object re-ID~\cite{qian2024reidtrends} with open-set identification~\cite{liao2014openset}.
The methods assessed in this study are SIFT~\cite{lowe2004sift} and ORB~\cite{rublee2011orb}, learned correspondences such as SuperGlue~\cite{sarlin2020superglue} and LoFTR~\cite{sun2021loftr}, and crop embeddings~\cite{zhou2019osnet,dosovitskiy2021vit,touvron2021deit,radford2021clip,oquab2024dinov2,jocher2023ultralytics}; all are applied zero-shot, without fine-tuning on \textsc{CrackID} (Section~\ref{sec:setup}).
LightGlue~\cite{lindenberger2023lightglue} and NetVLAD~\cite{arandjelovic2016netvlad} are left as future baselines (Section~\ref{sec:setup}).
Vehicle-damage re-ID across claims~\cite{maiano2023antifraud} follows our open-set protocol, whereas other studies focus on detection and severity rather than association~\cite{hasan2025review,perez2024assessment,mallios2023multiview}.
Aligning first and then reasoning about presence or absence is the approach of change detection~\cite{sakurada2015changedetection,alcantarilla2018street}, which relies on planar registration~\cite{hartley2004mvg,fischler1987ransac,barath2019magsac}; we add a fallback based on enhanced correlation coefficient (ECC) maximization~\cite{evangelidis2008ecc} (Section~\ref{sec:cascade}).
Our key distinction is that alignment and matching targets must be disjoint: the homography is estimated from the wall, not the crack.
